\part{TỔNG QUAN VÀ PHÂN TÍCH THỊ TRƯỜNG}

\chapter{Giới thiệu về ReactJS}

\section{Lịch sử hình thành và Bối cảnh ra đời}
Trước năm 2013, phần lớn các ứng dụng web trên thế giới được xây dựng dựa trên mô hình truyền thống: Server tạo ra các trang HTML tĩnh và gửi về Client. Khi các ứng dụng web ngày càng trở nên phức tạp, yêu cầu tương tác cao hơn, việc quản lý DOM (Document Object Model) thông qua Vanilla JavaScript hoặc jQuery dần bộc lộ nhiều điểm yếu. Mã nguồn trở nên rối rắm (Spaghetti code), khó bảo trì và đặc biệt là hiệu năng giảm sút nghiêm trọng khi có quá nhiều thao tác cập nhật giao diện cùng lúc.

Trong bối cảnh đó, ReactJS được phát triển bởi Jordan Walke - một kỹ sư phần mềm tại Facebook (nay là Meta). React ban đầu được tạo ra để giải quyết vấn đề cụ thể của Facebook: làm sao để cập nhật giao diện của tính năng News Feed và Chat một cách liên tục, theo thời gian thực mà không làm gián đoạn trải nghiệm người dùng. Năm 2013, ReactJS chính thức được phát hành dưới dạng mã nguồn mở và nhanh chóng tạo ra một cuộc cách mạng trong lĩnh vực lập trình Frontend.

\section{Kiến trúc và Nguyên lý hoạt động cốt lõi}
Sự thành công của ReactJS không đến từ việc cung cấp nhiều tính năng có sẵn, mà đến từ một triết lý thiết kế hoàn toàn mới. Các nguyên lý hoạt động cốt lõi bao gồm:

\begin{itemize}
    \item \textbf{DOM Ảo (Virtual DOM):} Đây là một trong những đột phá lớn nhất của React. Thay vì can thiệp trực tiếp vào DOM thật (một thao tác rất tốn kém tài nguyên), React tạo ra một bản sao nhẹ hơn của DOM trên bộ nhớ gọi là Virtual DOM. Khi trạng thái (State) thay đổi, React sẽ tạo ra một Virtual DOM mới, so sánh (diffing) nó với Virtual DOM cũ để tìm ra những điểm khác biệt tối thiểu. Sau đó, nó mới cập nhật chính xác những điểm khác biệt đó lên DOM thật (quá trình Reconciliation). Điều này giúp ứng dụng hoạt động cực kỳ mượt mà.

    \item \textbf{Kiến trúc hướng Component (Component-based Architecture):} React khuyến khích lập trình viên chia nhỏ giao diện người dùng thành các phần độc lập, có thể tái sử dụng gọi là Component. Ví dụ trong dự án Trello Mini, hệ thống được chia thành \texttt{Board}, \texttt{List}, \texttt{Card}. Mỗi Component tự quản lý trạng thái hiển thị của riêng nó, giúp mã nguồn dễ đọc, dễ kiểm thử và dễ bảo trì khi dự án phình to.

    \item \textbf{Luồng dữ liệu một chiều (One-way Data Flow):} Khác với cơ chế Binding hai chiều (Two-way Data Binding) của Angular, React áp dụng luồng dữ liệu một chiều chảy từ Component cha xuống Component con thông qua \texttt{props}. Điều này giúp việc theo dõi sự thay đổi của dữ liệu trở nên dễ đoán (predictable), giảm thiểu tối đa các lỗi liên quan đến trạng thái ứng dụng.
\end{itemize}

\section{Đánh giá và So sánh công nghệ (Critical Thinking)}
Để hiểu rõ tại sao nhóm chọn React cho đồ án này, cần có sự đánh giá và so sánh với các công nghệ khác:

\subsection{Tại sao không dùng Vanilla JS / jQuery?}
Với một ứng dụng như Kanban Board, việc kéo thả một Card từ List A sang List B đòi hỏi phải cập nhật lại giao diện của cả hai List, tính toán vị trí, và cập nhật mảng dữ liệu gốc. Nếu dùng JavaScript thuần, lập trình viên phải viết code tìm kiếm thủ công từng node DOM, gỡ bỏ và chèn node mới. Điều này rất dễ gây lỗi đồng bộ giữa UI và Dữ liệu. Với React, chúng ta theo triết lý \textit{State-Driven UI} (Giao diện được điều khiển bởi trạng thái): Lập trình viên không bao giờ chạm vào DOM, chỉ cần cập nhật mảng dữ liệu State, React sẽ tự động vẽ lại UI.

\subsection{ReactJS so với Vue.js và Angular}
\begin{itemize}
    \item \textbf{Angular:} Là một Framework đồ sộ, cung cấp sẵn mọi thứ (Routing, HTTP Client, Form Validation). Tuy nhiên, nó có \textit{learning curve} (đường cong học tập) rất dốc và kiến trúc đôi khi quá nặng nề cho các dự án vừa và nhỏ.
    \item \textbf{Vue.js:} Kết hợp được ưu điểm của cả React và Angular, dễ học, code trực quan. Tuy nhiên, hệ sinh thái và cộng đồng của Vue ở quy mô doanh nghiệp lớn vẫn chưa thể bằng React.
    \item \textbf{ReactJS:} Bản chất chỉ là một thư viện (Library) chứ không phải Framework. Nó cung cấp sự linh hoạt tối đa. Lập trình viên có quyền tự do lựa chọn thư viện đi kèm (như chọn \texttt{React Router} cho định tuyến, \texttt{Zustand} cho State Management). Điều này phù hợp với triết lý học tập: tự xây dựng từng viên gạch để hiểu rõ bản chất hệ thống.
\end{itemize}

\textbf{Phân tích Trade-off:} React cực kỳ phù hợp cho các Single Page App (SPA) mang tính tương tác phức tạp cao (như bảng Trello, Dashboard quản trị). Ngược lại, không nên dùng React cho các website dạng landing page tĩnh, nặng về SEO nội dung văn bản đơn thuần, vì nó sinh ra gánh nặng tải JavaScript không cần thiết ban đầu.

\chapter{Cơ hội nghề nghiệp và Nhu cầu thị trường}

\section{Vị trí công việc thực tế}
Việc thành thạo ReactJS mở ra rất nhiều cơ hội việc làm hấp dẫn và linh hoạt trong ngành Công nghệ Thông tin. Sinh viên sau khi trang bị kiến thức React có thể đảm nhận các vị trí:
\begin{itemize}
    \item \textbf{Frontend Developer / Web Engineer:} Chịu trách nhiệm thiết kế, xây dựng và tối ưu hóa giao diện tương tác phía người dùng (Client-side). Nhiệm vụ bao gồm chuyển đổi bản vẽ thiết kế (Figma, Sketch) thành mã nguồn thực tế, tối ưu hóa tốc độ tải trang (Core Web Vitals) và đảm bảo tính tương thích trên mọi thiết bị và trình duyệt (Responsive Web Design).
    \item \textbf{Fullstack Developer (MERN Stack / Next.js):} Kết hợp sức mạnh của ReactJS ở Frontend với các công nghệ Backend như NodeJS (Express/NestJS) hoặc hệ sinh thái Next.js để phát triển toàn bộ hệ thống từ Cơ sở dữ liệu, API đến Giao diện người dùng. Sự liền mạch trong việc sử dụng cùng một ngôn ngữ (JavaScript/TypeScript) giúp lập trình viên Fullstack tăng tốc độ phát triển dự án đáng kể.
    \item \textbf{Mobile App Developer (React Native):} Đây là một đặc quyền cực lớn của các lập trình viên React. Nhờ triết lý "Learn once, write anywhere" (Học một lần, viết mọi nơi), lập trình viên có thể tận dụng kiến thức về Component, Hooks, State Management để chuyển đổi sang viết ứng dụng di động cho cả hai nền tảng iOS và Android chỉ với một tập mã nguồn chung bằng \textbf{React Native}.
\end{itemize}

\section{Nhu cầu tuyển dụng trên thị trường}
Hệ sinh thái React không ngừng lớn mạnh và chứng minh được độ ổn định trong hơn một thập kỷ qua.
Theo các báo cáo và khảo sát thường niên từ các tổ chức uy tín thế giới như \textit{Stack Overflow Developer Survey} và \textit{State of JS}, ReactJS liên tục giữ vững ngôi vị là thư viện web được sử dụng phổ biến nhất và được cộng đồng lập trình viên yêu thích nhất.

Đặc biệt tại thị trường Việt Nam, dựa trên báo cáo phân tích xu hướng tuyển dụng từ các nền tảng việc làm IT hàng đầu như ITviec và TopCV, các từ khóa liên quan đến hệ sinh thái React ("React", "ReactJS", "NextJS", "React Native") luôn lọt top 3 kỹ năng được nhà tuyển dụng săn đón nhất. Sự tin tưởng này đến từ việc các tập đoàn công nghệ lớn trên thế giới (Meta, Netflix, Airbnb, Uber) cũng như các kỳ lân công nghệ tại Việt Nam (VNG, MoMo, Shopee, Tiki) đều đang xây dựng và vận hành các nền tảng có hàng triệu lượt truy cập mỗi ngày dựa trên lõi React.

\section{Mức lương tham khảo và Lộ trình thăng tiến}
Sự kết hợp giữa nhu cầu khổng lồ từ thị trường và yêu cầu chuyên môn cao khiến mức đãi ngộ cho lập trình viên ReactJS luôn duy trì ở mức rất cạnh tranh so với mặt bằng chung của ngành. Dưới đây là phân tích chi tiết về lộ trình thăng tiến và mức lương tham khảo tại thị trường Việt Nam (tùy thuộc vào năng lực thực tế, quy mô doanh nghiệp và khả năng ngoại ngữ):

\begin{itemize}
    \item \textbf{Fresher / Thực tập sinh (Mới ra trường - Dưới 1 năm kinh nghiệm):}
    Yêu cầu ở giai đoạn này là nắm vững cú pháp JavaScript (ES6+), hiểu nguyên lý hoạt động của Virtual DOM, vòng đời của Component (Lifecycle/Hooks) và thao tác gọi API cơ bản (Fetch/Axios). Công việc chủ yếu là code các component tĩnh và sửa các bug nhỏ gọn.
    Mức lương khởi điểm dao động từ \textbf{8.000.000 VNĐ đến 12.000.000 VNĐ / tháng}.

    \item \textbf{Junior / Mid-level (1 - 3 năm kinh nghiệm):}
    Ở cấp độ này, lập trình viên đã có khả năng hoạt động độc lập. Họ có thể xử lý các logic luồng dữ liệu phức tạp, am hiểu sâu về các kỹ thuật tối ưu hóa hiệu năng (hạn chế re-render bằng \texttt{useMemo}, \texttt{useCallback}), quản lý State thành thạo với các thư viện như Redux Toolkit, Zustand. Đồng thời có khả năng viết Unit Test cho UI.
    Mức lương dao động từ \textbf{15.000.000 VNĐ đến 28.000.000 VNĐ / tháng}.

    \item \textbf{Senior / Tech Lead (Trên 4 năm kinh nghiệm):}
    Vị trí này đòi hỏi tầm nhìn bao quát và tư duy kiến trúc (System Architecture). Họ là người ra quyết định công nghệ, thiết lập cấu trúc mã nguồn chuẩn mực (như Feature-Sliced Design), chịu trách nhiệm đánh giá mã nguồn (Code Review) và giải quyết các bài toán hóc búa nhất về hiệu năng hiển thị, bảo mật ứng dụng (XSS, CSRF) cho các hệ thống phức tạp tải lượng cao.
    Mức lương có thể vượt mức \textbf{40.000.000 VNĐ đến 60.000.000 VNĐ / tháng}, thậm chí cao hơn đáng kể đối với các cơ hội làm việc Remote cho các công ty ở Mỹ, Châu Âu, Úc hoặc Singapore.
\end{itemize}
